\documentclass[10pt,a4paper]{amsart}
\usepackage[margin=19mm]{geometry}
\usepackage{amsmath,amssymb,mathtools}
\usepackage[scr=rsfs,cal=euler]{mathalfa}
\usepackage{xfrac}
\usepackage[unicode,hidelinks,bookmarks=false]{hyperref}
\newcommand{\ie}{i.e.\ }
\newcommand{\cf}{cf.\ }
\newcommand{\resp}{respectively\ }
\newcommand{\Subsection}{\subsection}
\providecommand{\nobreakdash}{\nobreak\hbox{-}}
\setlength{\emergencystretch}{2em}
\multlinegap0pt
\usepackage{amssymb}
\usepackage{mathtools}
\newtheorem{definition}{Definition}
\newtheorem{theorem}{Theorem}
\title{On fixed point results in metric spaces for large triangle-perimeter contractions}
\author{Ovidiu Popescu}
\date{Source: arXiv:2605.13147v1 (CC BY 4.0)}
\begin{document}
\maketitle

\noindent\emph{Source and licence.} On fixed point results in metric spaces for large triangle-perimeter contractions. Version arXiv:2605.13147v1 (CC BY 4.0), distributed by arXiv under the Creative Commons Attribution 4.0 International licence, http://creativecommons.org/licenses/by/4.0/, as recorded in the arXiv licence metadata for this exact version and shown on the abstract page https://arxiv.org/abs/2605.13147v1. Full licence notice: \texttt{licenses/arxiv-2605.13147-CC-BY-4.0.txt}.\par\smallskip

\noindent\textit{Editorial scope.} Counted unit: the article's theorem theorem (source0.tex lines 216--218). The article's complete original proof of it is delivered with it (source0.tex lines 220--277, one \texttt{proof} environment of 58 lines). No mathematics is written, repaired or re-derived: the statement and the proof are the article's own lines, in the article's own notation, and no retained line is substituted or re-flowed.  Retained uncounted as the source-local context the proof needs: lines 146--156.  Nothing was omitted from any retained span, so the package carries no omission record.  The retained lines cite 1 works, whose entries are transcribed verbatim from the article's own bibliography source in the e-print and delivered with short editorial labels recorded in the item's reference mapping.  No retained line contains a figure environment, an image-inclusion command or drawing code, so no image is delivered and the package declares no asset dependency.  Editorial formatting only, in addition: an AMS \texttt{amsart} preamble that restates the article's own declarations and macros used by the retained lines, namely the environments \texttt{definition}, \texttt{theorem} on separate counters, together with the standard packages those lines need.  The retained environments print in this document as Definition 1, Theorem 1 in order of appearance, whereas the article prints the same items with the numbering of its own full document.  The complete native source of the article as distributed in the arXiv e-print for this exact version is bundled unchanged as \texttt{sources/200484/source0.tex}.
	\begin{definition}[Mesmouli \cite{Mesmouli}]
		Let $(X,d)$ be a metric space with  $|X|\geq 3$. A mapping $T \colon X\to X$ is a large triangle-perimeter contraction if 
		\begin{itemize}
			\item[(i)] for all pairwise distinct points $x,y,z \in X$, \(P(Tx,Ty,Tz) \leq \alpha P(x,y,z)\)
			\item[(ii)] for every \(\varepsilon > 0\), there exists a constant \(\delta(\varepsilon) \in (0,1)\) such that if
			\[\max\{d(x,y), d(y,z), d(x,z)\} \geq \varepsilon, \]
			then 
			\[P(Tx,Ty,Tz) \leq \delta(\varepsilon)  P(x,y,z). \]
		\end{itemize} 
	\label{largeP}
	\end{definition}

	\begin{theorem}
		Let $(X,d)$ be a complete metric space with  $|X|\geq 3$ and let $T \colon X\to X$ be a large triangle-perimeter contraction in the sense of Definition \ref{largeP}. Suppose \(T\) has no periodic points of prime period 2 and there exists $x_0 \in X$ and \(L>0\) such that \(d(x_0, T^nx_0)\leq L \), for all \(n \geq 1\). Then, \(T\) has a fixed point in \(X\). The number of fixed points is at most two.
	\end{theorem}

	\begin{proof}
		Let \(x_0 \in X\) such that \(d(x_0, T^nx_0)\leq L \), for all \(n \geq 1\). Consider the Picard iteration \((x_n)_{n\geq 0}\) defined by \(x_{n+1}:=Tx_n=T^{n+1}x_0\) and let for each \(n \geq 0\):
	\[P_n:= P(x_n, x_{n+1}, x_{n+2}) = d(x_n,x_{n+1})+d(x_{n+1},x_{n+2})+d(x_n,x_{n+2}).\]
		If there exists \(n\) such that $x_{n+1}=Tx_n$, then \(x_n\) is a fixed point. Otherwise, we have \(x_n\neq Tx_n\), for all \(n \geq 0\). Since $T$ has no periodic points of prime period \(2\), \(x_{n+2}=T^2x_n \neq x_n\), for each $n \geq 0$. Therefore, from condition (i) of Definition \ref{largeP}, we get 
	\[P_{n+1} = P(Tx_n, Tx_{n+1}, Tx_{n+2}) < P(x_n, x_{n+1}, x_{n+2}) = P_n,\]
		for every $n \geq 0$.
		
		Hence, \((P_n)\) is strictly decreasing and bounded below by \(0\).
		
		Suppose there exists positive integers \(m>n\), such that \(x_m=x_n\). Then, we have
		\[ x_{m+1} = Tx_m = Tx_n = x_{n+1}, \quad x_{m+2} = T^2x_m = T^2x_n = x_{n+2}. \]
		
		Thus, 
		\[ P_m = P(x_m, x_{m+1}, x_{m+2}) = P(x_n, x_{n+1}, x_{n+2}) = P_n, \]
		which contradicts the strict monotonicity of \((P_n)\). Therefore, if no fixed point occurs along the orbit, all \(x_n\) are distinct.
		
		Now, we prove that \((x_n)\) is a Cauchy sequence.
		
		Assuming that \((x_n)\) is not Cauchy, there exists \(\varepsilon_0 > 0\), such that, for every positive integer \(N\) we can find \(m > n > N\), with \(d(x_m,x_n) \geq \varepsilon_0\). Hence, there exist the sequences of indices \(m_k\), \(n_k\), with \(m_k>n_k \to \infty\), such that
		\[d(x_{m_k}, x_{n_k})\geq \varepsilon_0, \quad \text{for each } k \geq 0. \]
		
		By condition (i) of Definition \ref{largeP}, we get 
		\[P(x_{m_k+2}, x_{m_k+1}, x_{n_k+1}) < P(x_{m_k+1}, x_{m_k}, x_{n_k}) < P(x_{m_k}, x_{m_k-1}, x_{n_k-1}) < \dots \]
		\[\quad \qquad \qquad \qquad < P(x_{m_k-n_k+1}, x_{m_k-n_k}, x_{0}).\]
		
		Since \(d(x_{m_k}, x_{n_k}) \leq P(x_{m_k+1}, x_{m_k}, x_{n_k}),\) we get
		\[\varepsilon_0 \leq P(x_{m_k+1}, x_{m_k}, x_{n_k}) \leq P(x_{m_k}, x_{m_k-1}, x_{n_k-1}) \leq \dots \leq P(x_{m_k-n_k+1}, x_{m_k-n_k}, x_{0}).\]
		
		Since 
		\[\varepsilon_0 \leq P(a,b,c) \quad \text{implies} \quad \dfrac{\varepsilon_0}{3}\leq \max \{d(a,b), d(b,c), d(a,c)\}, \]
		we have \[ \dfrac{\varepsilon_0}{3}\leq \max \{d(x_{m_k+1-i},x_{m_k-i}), d(x_{m_k-i},x_{n_k-i}), d(x_{m_k+1-i},x_{n_k-i})\}, \]
		for each \(k \geq 0\) and \(i \in \{0,1,\dots, n_k\}\). Hence, by condition (ii) of Definition \ref{largeP}, we get
		\[ P(x_{m_k+1}, x_{m_k}, x_{n_k}) \leq \delta\left(\dfrac{\varepsilon_0}{3}\right) P(x_{m_k}, x_{m_k-1}, x_{n_k-1}) \leq \left[\delta\left(\dfrac{\varepsilon_0}{3}\right)\right]^2 P(x_{m_k-1}, x_{m_k-2}, x_{n_k-2})  \]
		\[\quad \qquad \qquad \qquad \leq \dots \leq \left[\delta\left(\dfrac{\varepsilon_0}{3}\right)\right]^{n_k}  P(x_{m_k-n_k+1}, x_{m_k-n_k}, x_{0}).\]
		
		However,
		
		\begin{align*} P(x_{m_k-n_k+1}, x_{m_k-n_k}, x_{0}) &= d(x_{m_k-n_k+1},x_{m_k-n_k})+ d(x_{m_k-n_k+1},x_0)+d(x_{m_k-n_k},x_0)\\ &\leq 2 d(x_{m_k-n_k+1},x_0) + 2 d(x_{m_k-n_k},x_0)\leq 4L, \end{align*}
		for each $k \geq 0$. Therefore, we obtain
		\[P(x_{m_k+1}, x_{m_k}, x_{n_k}) \leq 4L \cdot \left[\delta\left(\dfrac{\varepsilon_0}{3}\right)\right]^{n_k}.  \]
		Hence, we obtain,
		\[\varepsilon_0 \leq d(x_{m_k}, x_{n_k}) \leq P(x_{m_k+1}, x_{m_k}, x_{n_k}) \leq 4L \cdot \left[\delta\left(\dfrac{\varepsilon_0}{3}\right)\right]^{n_k} \to 0, \]
		as $k \to \infty$, which is a contradiction. Then, we conclude that \((x_n)\) is a Cauchy sequence.
		
		Since $(X,d)$ is complete and \((x_n)\) is Cauchy, there exists \(x^*\in X\) such that \(x_n\to x^*\) as \(n \to \infty\).
		
		If \(x_n \neq x^*\), for every \(n \geq 0,\) then, by condition (i) of Definition \ref{largeP}, we obtain 
		\[ P(x_{n+2}, x_{n+1}, Tx^*) < P(x_{n+1}, x_{n}, x^*) \to 0,\]
		as \(n \to \infty\), so \[d(x^*, Tx^*) = \lim\limits_{n\to \infty} d(x_{n+1}, Tx^*) \leq \lim\limits_{n\to \infty} P(x_{n+2}, x_{n+1}, Tx^*) = 0,\] i.e., \(x^*\) is a fixed point of \(T\). Otherwise, there exists \(n_0 \geq 0\) such that \(x^*=x_{n_0}\). Since all \(x_n\) are distinct, we have \(x_n \neq x^*\), for every \(n \geq n_0\) and by condition (i) of Definition \ref{largeP}, we get
		\[ P(x_{n+2}, x_{n+1}, Tx^*) < P(x_{n+1}, x_{n}, x^*) \to 0,\]
		as \(n \to \infty\), so \(d(x^*, Tx^*) = 0\), i.e., \(x^*\) is a fixed point of \(T\).
		
		Suppose \(x^*, y^*, z^* \in X\) are three distinct fixed points of \(T\). Then, by condition (i) of Definition \ref{largeP}, we get:
		\[P(x^*, y^*, z^*) = P(Tx^*, Ty^*, Tz^*) < P(x^*, y^*, z^*),\]
		which is a contradiction.
		
		Therefore, \(T\) has at most two fixed points.
	\end{proof}

\begin{thebibliography}{99}
\bibitem[Mesmou]{Mesmouli} M.B. Mesmouli, L.F. Iambor, T.S. Hassan, Fixed-Point Results in Metric Spaces for Large Triangle–Perimeter Contractions. Mathematics 2026, 14, 457.
\end{thebibliography}
\end{document}
